\subsection*{Transformer and main LV room proposal}
Propose two 1500 kVA transformers with separate LV sections:3000 kVA total gives 327.784 kVA margin. Allocate floors 1--5 to A and 6--10 to B, each with half the central plant/public allowance. Each section is 1336.108 kVA, about 2030 A at 380 V. The equal split is assumed. Full-load N--1 is not provided;  load shedding is required on loss of one unit. Any coupler must prevent parallel operation.

Place the rooms at the lower-left perimeter of the supplied plan: substation 12430 by 5450 mm and adjacent LV room 12430 by 3500 mm, the latter depth assumed. Reasons are external plant access, short LV interfaces, separation from the lift/stair/AHU core and independent LV access. Reserve a 3 m delivery route within 20 m of unloading, subject to external site confirmation.

The appendix shows assumed HV3000 by 1000 mm and transformer 2200 by 1800 mm equipment reservations and six LV cabinets;  these are not manufacturer footprints. Their clearance envelopes fit without overlap. COP101 v 16 clause 4.3 requires HV rear 1000/front 1800/sides 750 mm and transformer LV terminals 900/others 750 mm. Propose 2-hour partitions, sealed penetrations and FFL150 mm above external level. Outward 180-degree doors have shared sweeps requiring sequential use.

Two 600 by 2000 mm louvres at assumed 50\% free area provide 1.20 m$^2$, above clause 6.3's 1.12 m$^2$ reference;  certified free area must be at least 46.67\%. Fan symbols are conditional provisions, not actual heat-loss calculations. LV interfaces including assumed tails are 10.03/10.31 m, below the 20 m limit in 5.5.2.2. Actual equipment, structures, openings and ventilation must satisfy these declared conditions.
